\documentclass[10pt,a4paper]{amsart}
\usepackage{amsmath,amssymb,amsthm,graphicx}
% Retained preamble for the Stacks Project chapter snapshot in this directory.
% The chapter's own preamble.tex is a copy of the shared Stacks Project preamble; the block
% below is sliced MECHANICALLY, line for line, from it. Nothing was added, removed or
% redefined; the wrapper class, the xr-hyper/hyperref wiring and the externaldocument list
% are omitted because the assembler supplies the document class.
%
% lines taken from preamble.tex: 3, 5, 9, 11, 12, 156, 157, 158, ... (51 lines total)
%
% package coverage restored from the chapter's own preamble (2026-10-13): xy, xy-2cell

\usepackage[all]{xy}
\xyoption{2cell}
\UseAllTwocells

\newenvironment{reference}{\comment}{\endcomment}
\newenvironment{slogan}{\comment}{\endcomment}
\newenvironment{history}{\comment}{\endcomment}
\theoremstyle{plain}
\newtheorem{theorem}[subsection]{Theorem}
\newtheorem{proposition}[subsection]{Proposition}
\newtheorem{lemma}[subsection]{Lemma}
\theoremstyle{definition}
\newtheorem{definition}[subsection]{Definition}
\newtheorem{example}[subsection]{Example}
\newtheorem{exercise}[subsection]{Exercise}
\newtheorem{situation}[subsection]{Situation}
\theoremstyle{remark}
\newtheorem{remark}[subsection]{Remark}
\newtheorem{remarks}[subsection]{Remarks}
\numberwithin{equation}{subsection}
\def\lim{\mathop{\mathrm{lim}}\nolimits}
\def\colim{\mathop{\mathrm{colim}}\nolimits}
\def\Spec{\mathop{\mathrm{Spec}}}
\def\Hom{\mathop{\mathrm{Hom}}\nolimits}
\def\Ext{\mathop{\mathrm{Ext}}\nolimits}
\def\SheafHom{\mathop{\mathcal{H}\!\mathit{om}}\nolimits}
\def\SheafExt{\mathop{\mathcal{E}\!\mathit{xt}}\nolimits}
\def\Sch{\mathit{Sch}}
\def\Mor{\mathop{\mathrm{Mor}}\nolimits}
\def\Ob{\mathop{\mathrm{Ob}}\nolimits}
\def\Sh{\mathop{\mathit{Sh}}\nolimits}
\def\NL{\mathop{N\!L}\nolimits}
\def\CH{\mathop{\mathrm{CH}}\nolimits}
\def\proetale{{pro\text{-}\acute{e}tale}}
\def\etale{{\acute{e}tale}}
\def\QCoh{\mathit{QCoh}}
\def\Ker{\mathop{\mathrm{Ker}}}
\def\Im{\mathop{\mathrm{Im}}}
\def\Coker{\mathop{\mathrm{Coker}}}
\def\Coim{\mathop{\mathrm{Coim}}}
\DeclareMathSymbol{\boxtimes}{\mathbin}{AMSa}{"02}
\def\QCohstack{\mathcal{QC}\!\mathit{oh}}
\def\Cohstack{\mathcal{C}\!\mathit{oh}}
\def\Spacesstack{\mathcal{S}\!\mathit{paces}}
\def\Quotfunctor{\mathrm{Quot}}
\def\Hilbfunctor{\mathrm{Hilb}}
\def\Curvesstack{\mathcal{C}\!\mathit{urves}}
\def\Polarizedstack{\mathcal{P}\!\mathit{olarized}}
\def\Complexesstack{\mathcal{C}\!\mathit{omplexes}}
\def\Pic{\mathop{\mathrm{Pic}}\nolimits}
\def\Picardstack{\mathcal{P}\!\mathit{ic}}
\def\Picardfunctor{\mathrm{Pic}}
\def\Deformationcategory{\mathcal{D}\!\mathit{ef}}
\usepackage{cleveref}
\crefname{originallemma}{original lemma}{original lemmas}
\Crefname{originallemma}{Original lemma}{Original lemmas}
\makeatletter
\newlabel{equation-smooth}{{07XS}{source}{}{sourcepointer-1}{}}
\newlabel{equation-smooth@cref}{{[originallemma][0][]07XS}{[1][0][]source}{}{}{}}
\makeatother
\crefname{originallemma}{original lemma}{original lemmas}
\Crefname{originallemma}{Original lemma}{Original lemmas}
\makeatletter
\newlabel{formal-defos-lemma-versal-object-existence}{{06IW}{source}{}{sourcepointer-2}{}}
\newlabel{formal-defos-lemma-versal-object-existence@cref}{{[originallemma][0][]06IW}{[1][0][]source}{}{}{}}
\makeatother
\crefname{originallemma}{original lemma}{original lemmas}
\Crefname{originallemma}{Original lemma}{Original lemmas}
\makeatletter
\newlabel{lemma-approximate-versal}{{07XH}{source}{}{sourcepointer-3}{}}
\newlabel{lemma-approximate-versal@cref}{{[originallemma][0][]07XH}{[1][0][]source}{}{}{}}
\makeatother
\crefname{originallemma}{original lemma}{original lemmas}
\Crefname{originallemma}{Original lemma}{Original lemmas}
\makeatletter
\newlabel{lemma-deformation-category}{{07WU}{source}{}{sourcepointer-4}{}}
\newlabel{lemma-deformation-category@cref}{{[originallemma][0][]07WU}{[1][0][]source}{}{}{}}
\makeatother
\crefname{originallemma}{original lemma}{original lemmas}
\Crefname{originallemma}{Original lemma}{Original lemmas}
\makeatletter
\newlabel{lemma-implies-smooth}{{07XX}{source}{}{sourcepointer-5}{}}
\newlabel{lemma-implies-smooth@cref}{{[originallemma][0][]07XX}{[1][0][]source}{}{}{}}
\makeatother
\crefname{originallemma}{original lemma}{original lemmas}
\Crefname{originallemma}{Original lemma}{Original lemmas}
\makeatletter
\newlabel{lemma-versal-smooth}{{07XT}{source}{}{sourcepointer-6}{}}
\newlabel{lemma-versal-smooth@cref}{{[originallemma][0][]07XT}{[1][0][]source}{}{}{}}
\makeatother
\crefname{originallemma}{original lemma}{original lemmas}
\Crefname{originallemma}{Original lemma}{Original lemmas}
\makeatletter
\newlabel{remark-deformation-category-implies}{{07WV}{source}{}{sourcepointer-7}{}}
\newlabel{remark-deformation-category-implies@cref}{{[originallemma][0][]07WV}{[1][0][]source}{}{}{}}
\makeatother
\crefname{originallemma}{original lemma}{original lemmas}
\Crefname{originallemma}{Original lemma}{Original lemmas}
\makeatletter
\newlabel{section-axioms}{{07XJ}{source}{}{sourcepointer-8}{}}
\newlabel{section-axioms@cref}{{[originallemma][0][]07XJ}{[1][0][]source}{}{}{}}
\makeatother
\title{Producing a smooth chart through a versal finite-type family}
\author{The Stacks Project Authors}
\date{stacks-project:artin.tex:local-snapshot-20261012; GNU FDL 1.2 or later}
\begin{document}
\maketitle
\noindent\textit{Editorial note.} Statement span is lines 2407-2429 of sources/stacks-artin/pinned-original-source.tex; proof span is lines 2431-2454, the complete author proof (22 body lines).\par\medskip
\noindent\textit{Editorial note.} Navigation stubs: 8 cited-but-not-included labels are supplied as non-load-bearing pointers with their pinned tags where the snapshot tag index has them.\par\medskip
\noindent\textit{Editorial note.} extra\_packages is empty on purpose: the retained preamble already loads xy with options and a second declaration clashes in the assembler wrapper.\par\medskip
\noindent\textit{Editorial note.} Text is retained verbatim, including the author's own typos and placeholders.\par\medskip
\noindent\textit{Editorial note.} Queue index 36 of queue\_artin.json (pb=1090, stmt 2407-2429, proof 2431-2454).\par\medskip
\noindent\textit{Editorial note.} stacks-artin wave ar01 (2026-10-14): fresh adjudication accept per lane-s3/artin\_judgments.md; pool per Chief 10:5x direct dispatch (idle-first over S12's 01:5x note, SMC precedent); artin.tex 99 rows/67 with proof judged 46 accept / 21 reject / 5 AGP-overlap avoidance rows; patch-sd-envs contract spans scanned.\par\medskip
\noindent\textit{Allowed standard prerequisites (editorial).}\begin{itemize}
\item Explicit prerequisite boundary. The cited inputs are prerequisites; the argument in the retained span is not replaced by any cited result.
\end{itemize}\medskip
\ifdefined\hypertarget\hypertarget{sourcepointer-1}{}\fi
\noindent\textit{Source reference.} Original lemma 07XS: sources/stacks-artin/pinned-original-source.tex, source lines 2268--2275. Cited by the retained statement or proof but not part of the collected spans; supplied only as a navigation pointer. The argument does not depend on its content.\par\medskip
\ifdefined\hypertarget\hypertarget{sourcepointer-2}{}\fi
\noindent\textit{Source reference.} Original lemma 06IW: sources/stacks-artin/pinned-original-source.tex, source lines 2454--2454. Cited by the retained statement or proof but not part of the collected spans; supplied only as a navigation pointer. The argument does not depend on its content.\par\medskip
\ifdefined\hypertarget\hypertarget{sourcepointer-3}{}\fi
\noindent\textit{Source reference.} Original lemma 07XH: sources/stacks-artin/pinned-original-source.tex, source lines 2092--2102. Cited by the retained statement or proof but not part of the collected spans; supplied only as a navigation pointer. The argument does not depend on its content.\par\medskip
\ifdefined\hypertarget\hypertarget{sourcepointer-4}{}\fi
\noindent\textit{Source reference.} Original lemma 07WU: sources/stacks-artin/pinned-original-source.tex, source lines 824--833. Cited by the retained statement or proof but not part of the collected spans; supplied only as a navigation pointer. The argument does not depend on its content.\par\medskip
\ifdefined\hypertarget\hypertarget{sourcepointer-5}{}\fi
\noindent\textit{Source reference.} Original lemma 07XX: sources/stacks-artin/pinned-original-source.tex, source lines 2345--2354. Cited by the retained statement or proof but not part of the collected spans; supplied only as a navigation pointer. The argument does not depend on its content.\par\medskip
\ifdefined\hypertarget\hypertarget{sourcepointer-6}{}\fi
\noindent\textit{Source reference.} Original lemma 07XT: sources/stacks-artin/pinned-original-source.tex, source lines 2281--2288. Cited by the retained statement or proof but not part of the collected spans; supplied only as a navigation pointer. The argument does not depend on its content.\par\medskip
\ifdefined\hypertarget\hypertarget{sourcepointer-7}{}\fi
\noindent\textit{Source reference.} Original lemma 07WV: sources/stacks-artin/pinned-original-source.tex, source lines 866--885. Cited by the retained statement or proof but not part of the collected spans; supplied only as a navigation pointer. The argument does not depend on its content.\par\medskip
\ifdefined\hypertarget\hypertarget{sourcepointer-8}{}\fi
\noindent\textit{Source reference.} Original lemma 07XJ: sources/stacks-artin/pinned-original-source.tex, source lines 2462--2489. Cited by the retained statement or proof but not part of the collected spans; supplied only as a navigation pointer. The argument does not depend on its content.\par\medskip
\begin{lemma}
\label{lemma-get-smooth}
Let $S$ be a locally Noetherian scheme. Let
$p : \mathcal{X} \to (\Sch/S)_{fppf}$ be a category fibred in groupoids.
Let $k$ be a finite type field over $S$ and let $x_0$ be an object of
$\mathcal{X}$ over $\Spec(k)$ with image $s \in S$. Assume
\begin{enumerate}
\item $\Delta : \mathcal{X} \to \mathcal{X} \times \mathcal{X}$ is
representable by algebraic spaces,
\item $\mathcal{X}$ satisfies axioms [1], [2], [3] (see
Section \ref{section-axioms}),
\item every formal object of $\mathcal{X}$ is effective,
\item openness of versality holds for $\mathcal{X}$, and
\item $\mathcal{O}_{S, s}$ is a G-ring.
\end{enumerate}
Then there exist a morphism of finite type $U \to S$ and an object
$x$ of $\mathcal{X}$ over $U$ such that
$$
x : (\Sch/U)_{fppf} \longrightarrow \mathcal{X}
$$
is smooth and such that there exists a finite type point $u_0 \in U$
whose residue field is $k$ and such that $x|_{u_0} \cong x_0$.
\end{lemma}

\begin{proof}
By axiom [2], Lemma \ref{lemma-deformation-category}, and
Remark \ref{remark-deformation-category-implies}
we see that $\mathcal{F}_{\mathcal{X}, k, x_0}$ satisfies (S1) and (S2).
Since also the tangent space has finite dimension by axiom [3]
we deduce from Formal Deformation Theory, Lemma
\ref{formal-defos-lemma-versal-object-existence}
that $\mathcal{F}_{\mathcal{X}, k, x_0}$ has a versal formal object $\xi$.
Assumption (3) says $\xi$ is effective. By axiom [1] and
Lemma \ref{lemma-approximate-versal}
there exists a morphism of finite type $U \to S$, an object $x$ of
$\mathcal{X}$ over $U$, and a finite type point $u_0$ of $U$ with residue
field $k$ such that $x$ is versal at $u_0$ and such that
$x|_{\Spec(k)} \cong x_0$. By openness of versality we may shrink
$U$ and assume that $x$ is versal at every finite type point of $U$.
We claim that
$$
x : (\Sch/U)_{fppf} \longrightarrow \mathcal{X}
$$
is smooth which proves the lemma. Namely, by Lemma \ref{lemma-versal-smooth}
$x$ satisfies (\ref{equation-smooth})
whereupon Lemma \ref{lemma-implies-smooth}
finishes the proof.
\end{proof}

\end{document}
